% Reset dei contatori per l'appendice
\setcounter{figure}{0}
\renewcommand{\thefigure}{A.\arabic{figure}}
\setcounter{table}{0}
\renewcommand{\thetable}{A.\arabic{table}}

\section{Preprocessing del Target e Partitioning}
\label{app:data_splitting}

In questa sezione si descrivono le operazioni preliminari di codifica della variabile dipendente e la strategia di suddivisione del dataset (\textit{Data Splitting}) adottata per la fase di addestramento e validazione.

\subsection{Codifica della Variabile Target}

La variabile target originale \texttt{income}, di natura categorica nominale, è stata sottoposta a una procedura di pulizia e mappatura binaria per renderla idonea agli algoritmi di classificazione.
Al fine di garantire la robustezza del parsing contro eventuali artefatti di formattazione (e.g., spaziature o punteggiatura residua come \texttt{'<=50K.'}), è stata applicata una funzione di normalizzazione delle stringhe prima della codifica.
\noindent
La trasformazione è definita dalla seguente mappatura $M: \mathcal{Y}_{raw} \to \{0, 1\}$:
\begin{equation}
    y_i = 
    \begin{cases} 
    0 & \text{se } \texttt{income}_i \in \{\text{'<=50K'}\} \\
    1 & \text{se } \texttt{income}_i \in \{\text{'>50K'}\}
    \end{cases}
\end{equation}
La classe positiva ($y=1$) identifica gli individui con reddito superiore alla soglia, costituendo il target di interesse per la predizione.

\subsection{Strategia di Partitioning (Stratified Hold-out)}

Il dataset processato ($N=30.139$) è stato suddiviso in due sottoinsiemi disgiunti secondo un rapporto 80/20:
\begin{itemize}
    \item \textbf{Training Set ($S_{train}$):} 80\% dei dati ($n_{train} = 24.111$), utilizzato per la stima dei parametri del modello.
    \item \textbf{Test Set ($S_{test}$):} 20\% dei dati ($n_{test} = 6.028$), riservato esclusivamente per la valutazione delle performance su dati non osservati (\textit{unseen data}).
\end{itemize}

Per garantire la riproducibilità dell'esperimento, è stato fissato un \textit{seed} globale (\texttt{random\_state=42}).

\subsection{Campionamento Stratificato}
Dato il moderato sbilanciamento delle classi nel dataset originale (tasso di positività $\approx 24.9\%$), un campionamento casuale semplice (\textit{Simple Random Sampling}) avrebbe potuto introdurre bias campionari, alterando le probabilità a priori delle classi nei due sottoinsiemi.

\noindent
È stato pertanto adottato lo \textbf{Stratified Sampling}, vincolando lo split a preservare la proporzione originale delle classi:
\begin{equation}
    P(y=1 | S_{train}) \approx P(y=1 | S_{test}) \approx P(y=1 | S_{total})
\end{equation}

\noindent
La verifica a posteriori (Tabella \ref{tab:split_verify}) conferma l'efficacia della stratificazione, con una divergenza ($\Delta$) tra le distribuzioni del Training e del Test set inferiore allo $0.01\%$, garantendo che la valutazione sul test set sia rappresentativa della distribuzione reale.

\begin{table}[h]
\centering
\begin{tabular}{lccc}
\toprule
\textbf{Classe} & \textbf{Training Set \%} & \textbf{Test Set \%} & \textbf{Delta ($\Delta$)} \\
\midrule
0 ($\le$50K) & 75.0944\% & 75.0995\% & 0.0052\% \\
1 ($>$50K) & 24.9056\% & 24.9005\% & 0.0052\% \\
\bottomrule
\end{tabular}
\caption{Verifica della distribuzione delle classi post-split.}
\label{tab:split_verify}
\end{table}
